% Generado por ai_engine/figura_explicabilidad.py: no editar a mano.
\begin{figure}[ht]
\centering
\fbox{\parbox{0.93\textwidth}{\footnotesize\setlength{\baselineskip}{13pt}
Paciente de 70 años de edad, minero jubilado, sin alergias medicamentosas conocidas, que presenta como antecedentes personales: accidente laboral antiguo con fracturas vertebrales y costales; intervenido de enfermedad de Dupuytren en mano derecha y by-pass iliofemoral izquierdo; Diabetes Mellitus tipo II, \colorbox{amarillo!14}{hipercolesterolemia} e hiperuricemia; enolismo activo, fumador de 20 \colorbox{amarillo!45}{cigarrillos} / día. Es derivado desde Atención Primaria por presentar hematuria macroscópica postmiccional en una ocasión y microhematuria \colorbox{amarillo!14}{persistente} \colorbox{amarillo!28}{posteriormente,} con \colorbox{amarillo!28}{micciones} \colorbox{amarillo!14}{normales.} En la exploración física presenta un buen estado general, con abdomen y genitales \colorbox{amarillo!14}{normales;} tacto rectal compatible con
}}
\caption[Ejemplo de explicabilidad sobre un informe real]{Términos que sostienen la predicción
del código \texttt{F17.210} (Dependencia de nicotina, cigarrillos, sin complicaciones) sobre un informe del conjunto de
prueba. El sombreado es proporcional a la caída del \emph{logit} al enmascarar cada palabra,
en tres niveles. Los cinco primeros: \emph{cigarrillos} (0,41), \emph{posteriormente} (0,16), \emph{micciones} (0,14), \emph{persistente} (0,12), \emph{normales} (0,11).}
\label{fig:ejemplo-explicabilidad}
\end{figure}
